\documentclass[letterpaper, 12pt]{extarticle}

\usepackage[scale=.85]{geometry}

\linespread{2}

\usepackage{titlesec}
\titlelabel{(\thetitle) }

\usepackage{styles/common}%.sty
\usepackage{styles/ADDMC}%.sty
\usepackage{styles/docs}%.sty

\usepackage[bookmarksnumbered]{hyperref}

%%%%%%%%%%%%%%%%%%%%%%%%%%%%%%%%%%%%%%%%%%%%%%%%%%%%%%%%%%%%%%%%%%%%%%%%%%%%%%%%

\begin{document}

% \textbf{Latest version}: \vpHomepage \\ This version:
2019/11/10 \myskipbig

\textbf{STATEMENT OF RESEARCH} \\ \vp \myskipbig

%%%%%%%%%%%%%%%%%%%%%%%%%%%%%%%%%%%%%%%%%%%%%%%%%%%%%%%%%%%%%%%%%%%%%%%%%%%%%%%%

I am a \phd{} student in \cs{} at \rice{} (2017 - c. 2022).

%%%%%%%%%%%%%%%%%%%%%%%%%%%%%%%%%%%%%%%%%%%%%%%%%%%%%%%%%%%%%%%%%%%%%%%%%%%%%%%%
\myrule

\section{\pl}

In 2015/Summer, I started my first project in \pl{} under
the supervision of \nrlong.
This project was part of a NASA research contract with \kt,
\hi, and us.
\nr{} designed the logic programming language \llang{}
to encode NASA safety-critical procedures.
I implemented \llang{} by automating its translation into
\aplong{}
(another logic programming language
whose solver was already built).
In 2016/Spring, the project was completed.
The NASA \textit
{Disclosure of Invention and New Technology}
titled \disclosureNASA{}
recognized six innovators:
\nr, three other doctorate holders,
a \phd{} student, and me.

\myskip

In 2016/Spring, my second project in \pl{} with \nr{} was started.
He designed the \ledlong{} for literate programming
(inspired by the \code{WEB} system of \dknuth).
I implemented the executing aspect of \led{}
by automating its translation into the functional programming language \sllang.
To implement the formatting aspect,
I am developing an \led-to-\tex{} translator.

\newpage

In 2016/Fall, I started another project in \pl,
this time with \rglong.
This project is inspired by our
knot theory research back in 2015/Spring,
when we studied
the non-commutative A-ideal of the three-twist knot
by manually doing skein computations.
I want to automate these computations
by developing a computer algebra system.
So I designed the procedural programming language \kn.
Then I implemented \kn{}
using \code{SymPy} (a \python{} library)
for symbolic computation and \tex{} printing.
I am now implementing additional knot-theoretic operations
as \kn{} built-in functions.

\section{\ai}

In 2016/Summer, I started writing an \ai{} research paper
under the guidance of \mglong.
We investigate \cplong{},
which is a superset of the logic programming language \aplong{}
(whose stable model semantics was defined by \mg{} and \vl).
In detail, a \cp{} program is said to have
antichain property if
none of its answer set is a proper subset of another.
The problem is to find syntactic conditions
which guarantee this semantically desirable property.
We discovered such a sufficient condition which is
reasonably weak:
the program's dependency graph
being acyclic and
having no directed path from one cr-rule-head-vertex to
another.
I am revising the paper before submitting for publication.

%%%%%%%%%%%%%%%%%%%%%%%%%%%%%%%%%%%%%%%%%%%%%%%%%%%%%%%%%%%%%%%%%%%%%%%%%%%%%%%%

\end{document}
